\documentclass[]{article}
\usepackage[margin = .75in]{geometry}
\usepackage{hyperref}

\title{Intermediate Guide to LaTeX}
\author{GMA - UF Mathematics}
\date{A date or \today works}

\begin{document}
\maketitle
Some document classes have default parameters for \verb|title|, \verb|author|, \verb|date| which are set in the preamble by the commands: \verb|\title{title}|, \verb|\author{yours and collaborator names}|, \verb|\date{\today or a specific date}|. To then render these within the document, we use the command \verb|\maketitle|, which will render the above triplet. If \verb|\maketitle| is called after text, it will start a new page before rendering.

\section{Custom commands}
\verb|\newcommand|, \verb|\renewcommand|, \verb|arguments| (optional/required), \verb|\newenvironment|, 
\section{Counters}

\section{More details}
\subsection{Figures}
subcaption, wrapfig

\subsection{Tables}
Multi row/col
\subsection{Clever Refs}


\section{Version control}
\href{https://learngitbranching.js.org/}{Learn GIT}

\section{TikZ, pgfplots, and TikZ-cd}
\end{document}

**Workflow & Tooling**
- Understanding the full compilation pipeline (pdflatex → biber → pdflatex × 2)
- Makefiles and `latexmk` for automated builds
- Version control strategies for collaborative LaTeX projects
- Diffing LaTeX documents with `latexdiff`

**Debugging & Optimization**
- Interpreting and fixing common warnings (overfull/underfull hboxes)
- Suppressing vs. resolving package conflicts

The key shift at this level is moving from *using* LaTeX to *controlling* it — writing reusable commands, managing large projects, and producing output that looks genuinely professional rather than just functional.